\section{Discussion and limitations}
The most reproducible observation in this study is not a preferred training duration. It is that the gain added by fine-tuning depends strongly on the duration at which the model is evaluated. The symmetric intervention makes the distinction testable. Specialization predicts that moving the fine-tuning duration from $3.84$ to $10.24$\,s should move recovery toward the longer evaluation point. Instead, both fine-tunes gain more at $10.24$\,s and the long-trained checkpoint has the smaller duration interaction.

These results motivate a narrower hypothesis about \emph{expression of adaptation gains}. Fine-tuning may produce changes whose measurable benefit is amplified at some inference durations, even when those durations were not favored during adaptation. This account is consistent with the matched $20{,}000$-step intervention and with the public dense text-fine-tuned reference, which also shows a larger gain at $10.24$\,s. It is not equivalent to saying that short generation is broken. Under the primary CLAP endpoint, the dense model's own duration response is close to the real-audio response, so base-model degradation at $3.84$\,s does not explain the interaction. The mechanism that makes adaptation gains more visible at longer inference durations remains open.

The domain results illustrate why the same evaluation principle matters beyond duration. Clotho retains most of the AudioCaps recovery, whereas hip-hop retains only a small fraction of the dense gap despite a discriminating metric. A single label such as ``out of domain'' would hide this difference. Replications should therefore move along more than one domain axis and retain a dense or real-data anchor whenever a small gain could otherwise be mistaken for a metric floor.

The main causal limitation also remains clear. The experimental dense $2\times2$ uses a short raw-weight fine-tune against an EMA baseline and substantially degrades the dense model, so it cannot replace the unavailable dense checkpoint after the released $10^{6}$-step recovery. The matched $20{,}000$-step pruned intervention is strong evidence against training-duration specialization at that budget, but it does not determine whether the million-step duration dependence is specific to pruning, to fine-tuning in general, or to another property of AudioLDM.

Other limits concern generality and power. The study covers one model family, two pruning severities and a small domain grid. The severity~1 interaction varies substantially across two outcome-blind prompt samples. The step from $10.24$ to $15.36$\,s is too uncertain to establish a plateau. Human-CLAP, KL and PANNs corroborate the duration pattern but are still automatic metrics, while the small blinded author listening exercise is descriptive only.

The practical implication is simple. A recovery claim should report the paired increment from P to \PFT{} at the intended operating point and at deliberately displaced points. Common generation noise improves the comparison, and dense, chance or real-data anchors make the size of the gain interpretable. If a mechanism is proposed, its intervention should change a directional prediction rather than merely add another correlation. Here, matching the fine-tuning budget while changing training duration was the decisive test.

\section{Conclusion}
Recovery fine-tuning gain is strongly operating-point dependent in the AudioLDM-M checkpoints studied here. The duration interaction survives multiple scorers and both pruning severities, while a symmetric intervention shows that it does not follow training duration at a matched $20{,}000$-step budget. Recovery transfers strongly to Clotho and only weakly to hip-hop despite informative dense anchors. Reporting recovery as a paired function of operating point therefore gives a more faithful account of compressed-model behavior than a single post-fine-tuning score.

\clearpage
{\setlength{\itemsep}{0pt}\setlength{\parsep}{0pt}
\begin{thebibliography}{99}
\bibitem{fang2023diffpruning} G.~Fang, X.~Ma, and X.~Wang, ``Structural pruning for diffusion
models,'' in \emph{Proc.\ NeurIPS}, 2023.
\bibitem{kim2024bksdm} B.-K.~Kim, H.-K.~Song, T.~Castells \emph{et al.}, ``BK-SDM: A lightweight, fast,
and cheap version of Stable Diffusion,'' in \emph{Proc.\ ECCV}, 2024.
\bibitem{fang2025tinyfusion} G.~Fang, K.~Li, X.~Ma, and X.~Wang, ``TinyFusion: Diffusion transformers
learned shallow,'' in \emph{Proc.\ CVPR}, 2025.
\bibitem{liu2023audioldm} H.~Liu, Z.~Chen, Y.~Yuan \emph{et al.}, ``AudioLDM: Text-to-audio generation with latent diffusion models,'' in
\emph{Proc.\ ICML}, 2023.
\bibitem{singh2026pruning} A.~Singh, Y.~Yuan, Y.~Chen, W.~Wang, and M.~D.~Plumbley, ``Efficient
text-to-audio generation via pruning,'' arXiv:2607.13330, 2026.
\bibitem{wu2023clap} Y.~Wu, K.~Chen, T.~Zhang \emph{et al.},
``Large-scale contrastive language-audio pretraining with feature fusion and keyword-to-caption
augmentation,'' in \emph{Proc.\ ICASSP}, 2023.
\bibitem{huang2023makeanaudio} R.~Huang, J.~Huang, D.~Yang \emph{et al.}, ``Make-An-Audio:
Text-to-audio generation with prompt-enhanced diffusion models,'' in \emph{Proc.\ ICML}, 2023.
\bibitem{ghosal2023tango} D.~Ghosal, N.~Majumder, A.~Mehrish, and S.~Poria, ``Text-to-audio generation
using instruction guided latent diffusion model,'' in \emph{Proc.\ ACM Int.\ Conf.\ Multimedia (MM)},
2023, pp.~3590--3598.
\bibitem{kilgour2019fad} K.~Kilgour, M.~Zuluaga, D.~Roblek \emph{et al.}, ``Fr\'echet Audio
Distance: A reference-free metric for evaluating music enhancement algorithms,'' in
\emph{Proc.\ Interspeech}, 2019.
\bibitem{gui2024fad} A.~Gui, H.~Gamper, S.~Braun \emph{et al.}, ``Adapting Fr\'echet Audio
Distance for generative music evaluation,'' in \emph{Proc.\ ICASSP}, 2024.
\bibitem{chung2025kad} Y.~Chung, P.~Eu, J.~Lee \emph{et al.}, ``KAD: No more FAD!
An effective and efficient evaluation metric for audio generation,'' arXiv:2502.15602, 2025.
\bibitem{kong2020panns} Q.~Kong, Y.~Cao, T.~Iqbal \emph{et al.}, ``PANNs:
Large-scale pretrained audio neural networks for audio pattern recognition,'' \emph{IEEE/ACM
Trans.\ Audio, Speech, Lang.\ Process.}, vol.~28, pp.~2880--2894, 2020.
\bibitem{takano2025humanclap} T.~Takano, Y.~Okamoto, Y.~Kanamori \emph{et al.}, ``Human-CLAP:
Human-perception-based contrastive language-audio pretraining,'' in \emph{Proc.\ APSIPA ASC}, 2025,
pp.~131--136.
\bibitem{li2026finelap} X.~Li, X.~Xu, Z.~Ma \emph{et al.}, ``FineLAP: Taming
heterogeneous supervision for fine-grained language-audio pretraining,'' in \emph{Proc.\ ACL}, 2026,
pp.~10393--10408.
\bibitem{kumar2022finetuning} A.~Kumar, A.~Raghunathan, R.~Jones \emph{et al.}, ``Fine-tuning
can distort pretrained features and underperform out-of-distribution,'' in \emph{Proc.\ ICLR}, 2022.
\bibitem{kim2019audiocaps} C.~D.~Kim, B.~Kim, H.~Lee, and G.~Kim, ``AudioCaps: Generating captions
for audios in the wild,'' in \emph{Proc.\ NAACL-HLT}, 2019.
\bibitem{drossos2020clotho} K.~Drossos, S.~Lipping, and T.~Virtanen, ``Clotho: An audio captioning dataset,'' in \emph{Proc.\ ICASSP}, 2020.
\bibitem{agostinelli2023musiclm} A.~Agostinelli, T.~I.~Denk, Z.~Borsos \emph{et al.}, ``MusicLM:
Generating music from text,'' arXiv:2301.11325, 2023.
\bibitem{song2021ddim} J.~Song, C.~Meng, and S.~Ermon, ``Denoising diffusion implicit models,'' in
\emph{Proc.\ ICLR}, 2021.
\bibitem{wang2026ttabench} H.~Wang, C.~Liu, J.~Chen \emph{et al.}, ``TTA-Bench: A comprehensive benchmark for evaluating text-to-audio models,'' in \emph{Proc.\ AAAI}, 2026, pp.~33512--33520.
\bibitem{bibbo2026repo} G.~Bibb\'o, A.~Singh, and M.~D.~Plumbley, ``Companion repository for operating-point evaluation of recovery fine-tuning,'' 2026. [Online]. Available: \url{https://github.com/gbibbo/audioldm-modality-swap-pruning}
\end{thebibliography}}

\end{document}
